% !TeX root = ../main.tex
% ---------------------------------------------------------------------------
% Mechanised Tunnelling and TBM Engineering
% Scope extension: expands the TBM material sketched in Chapter 10 ( Chapter~\ref{ch:construction}).
% Key sources: Maidl et al., Mechanised Shield Tunnelling and Hardrock TBMs;
% Rostami & Ozdemir (CSM model); Bruland (NTNU model); Barton (QTBM);
% DAUB recommendations for TBM selection.
% ---------------------------------------------------------------------------
\chapter{Mechanised Tunnelling and TBM Engineering}
\label{ch:tbm-engineering}


\section{Machine selection}
\label{sec:machine-selection}

\subsection{Machine families and the selection decision}
\label{sec:tbm-families}
\scopenote{Open gripper, single shield, double shield, EPB, slurry and
multi-mode machines mapped against ground type, permeability, head and
alignment length. Selection is largely irreversible once procured.}

\subsection{Hard rock TBMs}
\label{sec:hard-rock-tbm}
\scopenote{Gripper machines and rock reaction, shield machines in variable
rock, and the support installation constraints imposed by a shielded
excavation.}

\section{Closed-face machines}
\label{sec:closed-face-machines}

\subsection{Earth pressure balance machines}
\label{sec:epb}
\scopenote{Chamber pressure control, screw conveyor discharge, plug stability,
and the conditioning required to convert granular soil into a plastic paste.}

\subsection{Slurry machines and separation}
\label{sec:slurry-tbm}
\scopenote{Bentonite circuit, air bubble pressure regulation, and the
separation plant as a governing constraint on advance rate in fine soils.}

\subsection{Multi-mode and variable-density machines}
\label{sec:multi-mode}
\scopenote{Convertible machines for alignments crossing several ground types;
conversion time and the cost of optionality.}

\section{Cutterhead and cutting tools}
\label{sec:cutterhead-and-cutting-tools}

\subsection{Cutterhead design}
\label{sec:cutterhead}
\scopenote{Cutterhead geometry, opening ratio, tool layout, spacing and
penetration interaction, and muck flow through the head.}

\subsection{Cutting tools and wear}
\label{sec:cutter-wear}
\scopenote{Disc cutters, drag picks and scrapers; abrasivity indices (CAI,
LCPC, Vickers); wear prediction; and cutter change as a programme and safety
driver.}

\section{Performance and installed power}
\label{sec:performance-and-installed-power}

\subsection{Penetration and advance rate prediction}
\label{sec:tbm-performance}
\scopenote{The CSM force-balance model, the NTNU empirical model and Barton's
$Q_{TBM}$. Utilisation as distinct from penetration rate, and why predicted
advance rates are so often optimistic.}

\subsection{Excavation energy and construction carbon}
\label{sec:tbm-energy-carbon}
\scopenote{The same ground--machine interaction that governs advance rate ---
torque, thrust, rate of penetration, rotation speed, diameter, depth, ground
class and in-situ stress --- also determines excavation energy, and through it
the A5 construction-stage carbon of Section~\ref{sec:carbon-tbm-energy}.
Machine performance is therefore a sustainability variable, not only a
programme variable.}

\subsection{Thrust, torque and installed power}
\label{sec:thrust-torque}
\scopenote{Thrust demand from cutting force, shield friction and back-up drag;
torque demand from cutterhead loading; sizing and the consequences of
under-specification.}

\section{Face pressure and annulus grouting}
\label{sec:face-pressure-and-annulus-grouting}

\subsection{Face pressure calculation and control}
\label{sec:tbm-face-pressure}
\scopenote{Setting the pressure envelope from the methods of
Chapter~\ref{ch:soft-ground}, converting it to machine operating parameters,
and controlling it in real time.}

\subsection{Annulus grouting and tail void closure}
\label{sec:annulus-grouting}
\scopenote{Simultaneous versus secondary grouting, grout mix and pressure
control, and its direct effect on the tail component of volume loss.}

\section{Guidance and interventions}
\label{sec:guidance-and-interventions}

\subsection{Guidance, steering and alignment control}
\label{sec:tbm-guidance}
\scopenote{Guidance systems, articulation, curve negotiation, roll control and
the build tolerance implications for the segmental ring.}

\subsection{Interventions and hyperbaric work}
\label{sec:hyperbaric}
\scopenote{Planned and emergency interventions, atmospheric versus compressed
air access, saturation diving techniques at high head, and the safety regime
that governs them.}

\section{Machine data and forward exploration}
\label{sec:machine-data-and-forward-exploration}

\subsection{Machine data and performance analysis}
\label{sec:tbm-data}
\scopenote{Continuous logging of thrust, torque, penetration, pressure and
grout volumes; using it for ground characterisation ahead of and behind the
face. Analytics in Chapter~\ref{ch:digital-underground}.}

\subsection{Probing and forward exploration}
\label{sec:tbm-probing}
\scopenote{Probe drilling, seismic-while-drilling and other look-ahead methods
from a shielded machine, and the programme cost of using them.}

\section{Risk and logistics}
\label{sec:risk-and-logistics}

\subsection{TBM risk: jamming, boulders, mixed face and karst}
\label{sec:tbm-risk}
\scopenote{The characteristic failure modes of mechanised tunnelling, the
conditions that produce them, and design and operational mitigations.}

\subsection{Launch, reception, and machine logistics}
\label{sec:tbm-logistics}
\scopenote{Launch shafts and cradles, seal design, reception, back-up
configuration, logistics and the surface site as a constraint on progress.}
